% =============================================================================
%  fault_location_id - manuscript_v2.tex
%
%  WP0.1-revised IEEE Access manuscript for the HIF transfer-function locator.
%  Closes deliverable D-A of the v3 execution plan
%  (docs/FaultLocationIdentification_ExecutionPlan.pdf §3, §6).
%
%  Authoring conventions
%  ---------------------
%  Every headline numerical value the abstract / §VI / §IX / glossary
%  must agree on is defined exactly once via a \newcommand macro below.
%  Reuse the macro everywhere; never hand-type the value.  This is the
%  mechanism that satisfies the WP0.1 acceptance criterion
%      "numerical values for location error, Rx error and CRLB headline
%       are byte-identical wherever they appear in the manuscript."
% =============================================================================

\documentclass[journal]{IEEEtran}

\usepackage[T1]{fontenc}
\usepackage[utf8]{inputenc}
\usepackage{amsmath,amssymb,amsfonts}
\usepackage{graphicx}
\usepackage{booktabs}
\usepackage{siunitx}
\sisetup{detect-all}
\usepackage{cite}
\usepackage{url}
\usepackage{hyperref}
\hypersetup{colorlinks=true,linkcolor=blue,citecolor=blue,urlcolor=blue}

% --------------------------------------------------------------------------
%  Headline numerical macros (single source of truth) - see WP0.1 spec
% --------------------------------------------------------------------------
\newcommand{\headlineLocErrNoiseless}{0.009\,\%}
\newcommand{\headlineLocErrSNR}%
  {1.18\,\% at $\mathrm{SNR}_I = 20$\,dB}
\newcommand{\headlineRxErrNoiseless}{2.42\,\%}
\newcommand{\headlineRxErrEnvelope}%
  {$\leq 8$\,\% across 100\,--\,2000\,$\Omega$ at SNR $\geq 30$\,dB}
\newcommand{\headlineCRLBGap}%
  {within 15\,\% of the Cram\'er--Rao Lower Bound at SNR $\geq 40$\,dB}
\newcommand{\samplingConfig}%
  {$F_s = 10$\,kHz, $N_s = 200$, one cycle window}

% -- Phase-2 macros (WP2.6 source-of-truth additions) ----------------------
% K03 = distributed-vs-50sec forward-model accuracy (P2.3 acceptance).
% K04 = cross-platform Phase-2 vs Phase-1 estimator improvement (P2.5);
%       the negative sign is real and reported honestly; the diagnosis
%       and forward-pointer to WP3.5 / WP3.6 are the subject of \S\,VI-D.
\newcommand{\headlineKThreeAcc}%
  {$\bar{\varepsilon}_H = 4.3\!\times\!10^{-6}\,\%$
   (max $2.7\!\times\!10^{-5}\,\%$)}
\newcommand{\headlineKFourImp}%
  {$-830.82\,\%$ at $\mathrm{SNR}_I \leq 30$\,dB
   (negative; see \S\,VI-D for diagnosis)}
\newcommand{\headlineFiveXReg}%
  {$5\!\times$ improvement vs the v1 R-L-only ceiling on the worst
   $(\alpha, R_x) = (0.95, 5000)$ cell}

% Glossary cross-reference helper
\newcommand{\glossref}[1]{\textsuperscript{(see Glossary, \texttt{#1})}}

% =============================================================================
\begin{document}
% =============================================================================

\title{Single-Ended Joint Estimation of HIF Location and Arc Resistance
       via Power-Frequency Admittance Identification with
       Dual-Channel Noise Modelling}

\author{%
  Arjundas K. and K. Shanthi Swarup,~\IEEEmembership{Senior Member,~IEEE}%
  \thanks{The authors are with the Power Systems Computational Lab,
          Department of Electrical Engineering,
          Indian Institute of Technology Madras, Chennai 600\,036, India
          (e-mail: arjundas@iitm.ac.in; swarup@ee.iitm.ac.in).}%
  \thanks{This work was conducted under the SAMBPS Digital Twin Labs
          R\&D programme on Smart Adaptive Model-Based Protection \&
          Control Systems (SAMBPS DTaaS).  The authors gratefully
          acknowledge the SAMBPS DTaaS reproducibility infrastructure
          and the public release of all 720-case waveform data,
          50-section reference dataset and analysis scripts under MIT
          licence.}%
}

\markboth{IEEE Access (submitted 2026)}{Arjundas \& Swarup: HIF-TF Locator (v2)}

\maketitle

% =============================================================================
\begin{abstract}
% (a) motivation - one sentence
High-impedance faults (HIFs) account for an estimated 5--10\,\% of
distribution-feeder faults and roughly a quarter of downed-conductor
events go undetected by conventional overcurrent protection
(IEEE PSRC Working Group D15, 1996), and grid-caused ignitions
represent 19\,\% of all U.S. wildfires over 2016--2020
(NREL TP-5R00-80746, 2023), making accurate HIF location and arc
resistance estimation an urgent safety problem.
% (b) method line
This paper proposes a single-ended joint estimator of fault location
$\alpha$ and arc resistance $R_x$ that fits a state-space transfer
function $H_{\mathrm{model}}(j\omega_0;\alpha,R_x)$ to the
power-frequency single-bin DFT admittance
$H_{\mathrm{meas}}(j\omega_0)$ via a two-stage optimiser
(grid + multi-start, then gradient with Armijo line-search) over a
720-case sweep ($10\alpha\times 5R_x\times 4\,\mathrm{SNR}_V\times
4\,\mathrm{SNR}_I$) under independent dual-channel additive white
Gaussian noise.
% (c) headline numbers
The estimator achieves a mean fault-location error of
\headlineLocErrNoiseless{} in the noiseless baseline and
\headlineLocErrSNR{}\glossref{LOC-ERR}; mean arc-resistance error is
\headlineRxErrNoiseless{} noiseless and \headlineRxErrEnvelope%
\glossref{RX-ERR}; the location-error metric is distinct from the
TF-magnitude modelling error of the two-section $\pi$-model versus a
50-section reference, which is reported separately in \S\,II.D.
% (d) CRLB headline
The empirical RMS estimator stays \headlineCRLBGap{}\glossref{CRLB-GAP}.
% (d') Phase-2 forward-model upgrade headline
A continuously parametrised distributed-parameter forward model
(cascaded ABCD blocks) replaces the two-section $\pi$ stack of v1
and reproduces a 50-section $\pi$ reference to \headlineKThreeAcc{};
on the worst $(\alpha, R_x)$ cell of \S\,II-D this is
\headlineFiveXReg.  Cross-platform empirical estimator improvement
from this upgrade is positive on self-consistent waveforms but
saturates against the single-bin DFT identifiability floor on
PSCAD / EMTP / 50-section reference data --- a finding we report
honestly and trace to a corrected-CRLB structural limit (\S\,VIII)
that resolves at the multi-bin / multi-port observation extension
deferred to a planned IEEE TPWRD follow-on.
% (e) sampling configuration
All results use \samplingConfig{}\glossref{SAMP-CFG}.
% (f) outlook
Single-frequency operation, training-free deployment, and the absence
of remote-end communication make the estimator a practical
single-ended primary HIF locator suitable for integration into the
SAMBPS DTaaS Protection-Validation digital-twin module.
\end{abstract}

\begin{IEEEkeywords}
High-impedance fault, fault location, arc resistance, transfer-function
identification, single-ended, joint estimation, training-free,
single-frequency, Cram\'er--Rao Lower Bound, dual-channel noise model,
distribution protection, SAMBPS DTaaS.
\end{IEEEkeywords}

% =============================================================================
\section{Introduction}\label{sec:intro}
% =============================================================================

% ---- Wildfire / safety motivation (WP0.2 sub-task 2) ---------------------
High-impedance faults (HIFs) account for an estimated 5--10\,\% of all
distribution-feeder faults, and approximately one quarter of
downed-conductor incidents go undetected by conventional overcurrent
relaying~\cite{PSRC1996D15}.  In the United States, grid-related
ignitions caused 19\,\% of wildfires recorded between 2016 and
2020~\cite{NREL2023TP5R0080746}.  Two catastrophic events frame the
operational stakes: the 2009 Australian Black Saturday bushfires
killed 173 people, with electrical assets implicated as ignition
sources on multiple fire fronts~\cite{BlackSaturday2009RoyalCommission};
the 2018 PG\&E Camp Fire killed 85 and led to a USD 13.5 billion
settlement after a worn transmission hook ignited dry grassland near
Pulga, California~\cite{CampFire2018PGE}.

In response, the California Public Utilities Commission (CPUC)
implemented Senate Bill 901 (SB\,901), placing utilities under a
formal Public Safety Power Shutoff (PSPS) framework that mandates
de-energisation when fire-weather risk exceeds defined
thresholds~\cite{CPUC2018SB901}.  PSPS shutoffs avoid catastrophic
ignition but at significant social and economic cost; selective HIF
location and arc-resistance estimation, by contrast, would enable
targeted de-energisation of only the faulted segment, materially
shrinking the PSPS footprint.  The estimator developed in this paper
is designed to fill that gap.

% ---- Prior-art classification (WP0.2 sub-task 1) -------------------------
The prior-art literature on HIF location and detection clusters into
seven distinct families, summarised in Table~\ref{tab:taxonomy}.  Two
columns are added for the comparison this paper draws:
``Single-ended,'' which marks methods that operate from a single line
terminal, and ``Joint $\alpha\!+\!R_x$,'' which marks methods that
estimate fault location $\alpha$ and arc resistance $R_x$
\emph{simultaneously} from the same observation rather than estimating
one and treating the other as a fixed parameter or post-processing
step.

\begin{table*}[t]
\caption{Seven-family taxonomy of prior-art HIF location / detection
methods.  ``Single-ended'' marks methods that use only one line
terminal; ``Joint $\alpha\!+\!R_x$'' marks methods that estimate both
quantities simultaneously from a single observation.  None of the
seven families is Yes in both columns.}
\label{tab:taxonomy}
\centering
\renewcommand{\arraystretch}{1.18}
\begin{tabular}{@{}p{2.6cm}p{6.6cm}cc p{3.4cm}@{}}
\toprule
\textbf{Family} & \textbf{Representative anchors} &
\textbf{Single-ended} & \textbf{Joint $\alpha\!+\!R_x$} &
\textbf{Distinguishing trait}\\
\midrule
(i) Signal processing / morphology / wavelet
  & Aucoin--Russell 1987 IEEE TPWRD~\cite{AucoinRussell1987TPWRD};
    Brahma 2013 IEEE TPS~\cite{Brahma2013TPS};
    Mishra 2019 IEEE TII (VMD)~\cite{Mishra2019TII}
  & Yes & No
  & Detection-only; no $R_x$ output.\\
(ii) Impedance / admittance / spectral
  & Iurinic--Bretas 2018 IJEPES~\cite{IurinicBretas2018IJEPES};
    Orozco-Henao--Bretas 2020 EPSR~\cite{OrozcoHenaoBretas2020EPSR};
    Penaloza 2023 EPSR~\cite{Penaloza2023EPSR}
  & Yes & No
  & $R_x$ enters as a parameter, not jointly estimated.\\
(iii) Two-ended impedance / synchrophasor
  & Eriksson--Saha 1985 IEEE TPAS~\cite{ErikssonSaha1985TPAS};
    Zeng 2021 EPSR~\cite{Zeng2021EPSR};
    Chandran 2024 IET GTD~\cite{Chandran2024IETGTD}
  & No & No
  & Requires comm channel + remote-end measurement.\\
(iv) Traveling-wave
  & SEL T401L~\cite{SEL_T401L_TM};
    Costa--Lopes 2015--2022~\cite{CostaLopes2015};
    IEEE TPWRD 2024~\cite{IEEE_TPWRD_TW_2024}
  & Yes & No
  & MHz-bandwidth front; HIF wavefront weak / smeared.\\
(v) $\mu$-PMU / data-driven
  & Cui--Weng 2020 IEEE TSG~\cite{CuiWeng2020TSG};
    Shahsavari 2019 IEEE TSG~\cite{Shahsavari2019TSG};
    Farajollahi 2018~\cite{Farajollahi2018}
  & No & No
  & Requires distributed $\mu$-PMU deployment.\\
(vi) ML / DL / GAN / GNN
  & Sirojan 2022 NCAA (transformer-CNN)~\cite{Sirojan2022NCAA};
    Veerasamy 2021 IEEE Access (LSTM)~\cite{Veerasamy2021Access};
    RGATv2 arXiv 2510.03571 (2025)~\cite{RGATv2_2025_arXiv}
  & Yes & No
  & Training-data dependent; opaque to operator.\\
(vii) Eigenvalue / state-matrix
  & Paramo--Bretas 2023 ISGT~\cite{ParamoBretas2023ISGT};
    Bretas 2021 MDPI Appl.\ Sci.~\cite{Bretas2021MDPI}
  & Yes & No
  & Recovers eigenstructure; $R_x$ implicit only.\\
\bottomrule
\end{tabular}
\end{table*}

\textit{The proposed estimator described in this paper would occupy an
eighth row of Table~\ref{tab:taxonomy} and is the only entry that is
``Yes'' in both the Single-ended and the Joint $\alpha\!+\!R_x$
columns: it operates from a single line terminal and produces a
simultaneous estimate $(\hat\alpha,\hat R_x)$ from a single complex
admittance observation $H_{\mathrm{meas}}(j\omega_0)$.}

Of the seven prior-art families, only (i), (ii), (iv), (vi) and (vii)
operate from a single terminal, and none produces a simultaneous
estimate of fault location and arc resistance.  Families that do
report $R_x$ either treat it as a nuisance parameter solved for after
location is fixed (family (ii)) or infer it implicitly through a
classifier boundary (family (vi)).

% ---- Contributions (WP0.2 sub-task 3 - C5 reframed) ----------------------
\subsection*{Paper contributions}
\begin{enumerate}
  \item A single-ended power-frequency admittance identification
        formulation that derives $H_{\mathrm{meas}}(j\omega_0)$ from
        the dual-channel single-bin DFT and matches it against a
        state-space transfer function $H_{\mathrm{model}}(j\omega_0;
        \alpha, R_x)$ derived in closed form from the two-section
        $\pi$-line.
  \item A two-stage optimiser (grid + multi-start, then gradient
        descent with Armijo line-search) that jointly estimates
        $(\hat\alpha,\hat R_x)$ from a single complex observation,
        without remote-end communication and without offline training.
  \item A 720-case study under independent dual-channel additive
        white Gaussian noise that reports per-cell mean and
        95\textsuperscript{th}-percentile location and $R_x$ error,
        including a noiseless baseline that isolates the structural
        floor.
  \item An explicit disambiguation of two error metrics that the
        prior literature routinely conflates:
        (a) the TF-magnitude modelling error of the two-section
        $\pi$-model versus a 50-section reference, and
        (b) the location error of the estimator
        (Glossary entries \texttt{MODEL-ERR} and \texttt{LOC-ERR},
        \S\ref{sec:results}).
  \item An \emph{identifiability bound}, derived from the
        Cram\'er--Rao inequality, that quantifies the lowest variance
        any unbiased single-bin admittance estimator can achieve at a
        given $(\alpha, R_x, \mathrm{SNR}_V, \mathrm{SNR}_I)$.  We
        position this not as a methodological novelty but as the
        objective reference against which the proposed estimator's
        residual gap is measured (\headlineCRLBGap).
\end{enumerate}

% ---- Forward-reference roadmap (WP0.2 sub-task 4) ------------------------
\noindent\textbf{Roadmap.}
Section~\ref{sec:models} sets up the distribution-line and arc
\emph{modelling} (two-section $\pi$-line, anti-parallel diode arc, and
the section-model accuracy study);
Section~\ref{sec:tf} presents the transfer-function
\emph{identification} (single-bin DFT cost and analytical model
gradient);
Section~\ref{sec:opt} describes the two-stage \emph{optimiser};
Sections~\ref{sec:sim}--\ref{sec:results} carry the
\emph{validation} programme (simulation matrix, dual-channel noise
model, results);
Section~\ref{sec:compare} compares against prior-art numerically; and
Section~\ref{sec:crlb} derives the Cram\'er--Rao identifiability
\emph{bound} that the estimator is benchmarked against.
Section~\ref{sec:conclusion} concludes with a 12\,/\,24\,/\,36-month
roadmap into the SAMBPS DTaaS Protection-Validation digital-twin
module.

\section{System and Models}\label{sec:models}

\subsection{Distribution-line parameterisation (\S\,II.A)}\label{sec:lineparam}
The faulted feeder is an \SI{11}{kV}, \SI{100}{km} radial overhead
line with per-unit-length parameters
$R' = \SI{0.0728}{\ohm/km}$,
$L' = \SI{0.927}{mH/km}$,
$C' = \SI{11.6}{nF/km}$
(Saha~2010, Springer Table 3.1~\cite{Saha2010BookFL}).
The high-impedance fault appears as a single shunt resistance $R_x$
at per-unit position $\alpha\in(0,1)$ measured from the source
substation.  All identification is performed at the
power-frequency carrier $f_0 = 50$\,Hz and a single complex DFT bin.

\subsection{Continuously parametrised forward model (\S\,II.B)}
\label{sec:fwdmodel}
The optimiser of \S\,\ref{sec:opt} couples to a
\emph{continuously parametrised distributed-parameter} forward model
$H(j\omega_0;\alpha,R_x)$ landed at WP2.1.  Cascaded ABCD blocks
with propagation constant $\gamma = \sqrt{(R'+j\omega_0 L')(G' +
j\omega_0 C')}$ and characteristic impedance
$Z_c = \sqrt{(R'+j\omega_0 L')/(G' + j\omega_0 C')}$ give
\begin{equation}
\!\!T(L) = \begin{bmatrix}
   \cosh(\gamma L) & Z_c\,\sinh(\gamma L)\\
   \sinh(\gamma L)/Z_c & \cosh(\gamma L)
\end{bmatrix},\quad
T_f = \begin{bmatrix}1 & 0\\ 1/R_x & 1\end{bmatrix},
\label{eq:abcd}
\end{equation}
and
$T = T_1(\alpha L)\,T_f\,T_2((1-\alpha)L)\,T_{\text{load}}$ with
$T_{\text{load}} = \mathrm{diag}(1, R_{\text{load}})$ ($R_{\text{load}}
= \SI{1}{\mega\ohm}$, open far-end consistent with the §II-A
deployment assumption).  The sending-end admittance follows by
inversion of the resulting two-port:
$H = C_{\text{end}}/A_{\text{end}}$.  The full derivation is in
Appendix~A~\S\,A.4--A.5 and \S\,A.7; the implementation is in
\texttt{models/faultloc\_distributed\_param\_model.py}.

The cascaded-$\Gamma$ two-section $\pi$-model state-space of P0.5
(implementation in \texttt{models/faultloc\_pi\_section\_model.py},
Appendix~A~\S\,A.3) is retained as a \emph{baseline demonstrator}
only: \S\,II-D quantifies that its modelling-error envelope
($\bar{\varepsilon}_H = 0.39\,\%$) is more than four orders of
magnitude looser than the closed-form distributed model's
\headlineKThreeAcc{}, but it remains useful for analytical
clarity (the $4\!\times\!4$ state-space exposes the
$A_{11} = -1/(R_x C_1)$ structure that drives the cost-surface
ill-conditioning analysed in \S\,\ref{sec:crlb}).  Phase 2 onwards
(WP2.4, this paper) the optimiser invokes (\ref{eq:abcd}) by default;
the cascaded-$\Gamma$ stack is exercised only in the legacy-comparison
test path.

\subsection{Anti-parallel diode arc model (\S\,II.C)}
\label{sec:arcmodel}
The HIF arc is represented by an anti-parallel diode pair with
asymmetric forward-conduction characteristics following the
Aucoin--Russell-1987 trace family~\cite{AucoinRussell1987TPWRD}.  The
PSCAD~\cite{Penaloza2023EPSR} and EMTP-RV cases of \S\,V exercise
this arc model with provenance noted; arc-class diversification
(Cassie--Mayr--Kizilcay; Wang-2020 distortion-controllable;
Torres-2022 stochastic) is the subject of WP4.4 in the Phase-4
roadmap (\S\,\ref{sec:conclusion}).

\subsection{Section-model accuracy study (\S\,II.D)}\label{sec:modeldelta}

The convergence study below has been updated under WP2.3 to include
the closed-form distributed-parameter model from §IV.  Three forward
models are evaluated against a 50-section $\pi$ reference (Saha-2010
half-$\pi$ convention) on the 19$\,\alpha\!\times\!5R_x$ operating
grid at $f_0 = 50$\,Hz; magnitude errors $\bar{\varepsilon}_H$ in
percent.

\begin{table}[!t]
\caption{Forward-model accuracy vs the WP1.3 50-section $\pi$
reference on the operating grid.  ``Mean'' is across all 95 cells;
``Max'' is the worst cell.  Updated under WP2.3 (P2.3).}
\label{tab:modelfit}
\centering
\renewcommand{\arraystretch}{1.18}
\begin{tabular}{@{}lrrl@{}}
\toprule
\textbf{Forward model} & \textbf{Mean $|\Delta H|/|H_{\text{ref}}|$} &
\textbf{Max} & \textbf{Source}\\
\midrule
v1 R-L-only 2-section (no shunt $C$) & $40.4\,\%$ & $87.5\,\%$
   & \texttt{models/faultloc\_legacy\_v1\_2section.py}\\
P0.5 cascaded-$\Gamma$ 2-section
   & $0.39\,\%$ & $0.98\,\%$
   & \texttt{models/faultloc\_pi\_section\_model.py}\\
\textbf{P2.1 distributed-parameter (closed form)}
   & $\mathbf{4.3\!\times\!10^{-6}\,\%}$ & $\mathbf{2.7\!\times\!10^{-5}\,\%}$
   & \texttt{models/faultloc\_distributed\_param\_model.py}\\
\bottomrule
\end{tabular}
\end{table}

\textbf{Headline:} the closed-form distributed-parameter model
agrees with the 50-section $\pi$ reference to better than
$\mathbf{1\,\%}$ across the operating grid (D-C acceptance).
Empirically the agreement is at numerical precision, set by the
50-section reference's own sectioning floor (which we measured at
$\sim 2\!\times\!10^{-5}\,\%$ vs a 100-section reference --- the
residual is the WP1.3 lumped-line discretisation, not a real
distributed-vs-$\pi$ gap).  The per-cell breakdown is in
\texttt{outputs/phase2\_reproduction.csv} (95 rows; \texttt{source\_of\_residual}
column = \texttt{within\_target} on every row).

\textbf{Out of scope.}  Frequency-dependent line parameters
($R'(\omega)$ skin effect; earth-return $R_g'(\omega)$) are present
in PSCAD's J.\,Marti FD line and EMTP-RV's FD model, but neither the
P2.1 closed-form distributed model nor the WP1.3 50-section $\pi$
reference includes them.  The Phase-1 cross-platform validation
(P1.4) showed that this lumped-vs-FD gap, at $f_0 = 50$\,Hz with the
short-feeder geometry studied here, is $<\!1\,\%$ in $|H|$ across
the grid (P1.3 finding); WP4.1 will quantify the impairment-class
sensitivity.  Frequency-dependent series impedance is therefore
deferred to Phase~4 and treated as an \emph{intentional} omission
at the §IV / §V layer.

\section{Transfer-Function Identification}\label{sec:tf}

The single-ended measurement is the complex sending-end admittance
at the power-frequency carrier.  Sampling the discrete-time
voltage $v[n]$ and current $i[n]$ at $F_s = 10$\,kHz over one
50\,Hz cycle ($N_s = 200$ samples) and taking the single-bin DFT
at $f_0$ gives
\begin{equation}
V_{\mathrm{bin}} = \sum_{n=0}^{N_s-1} v[n] \, e^{-j 2\pi n / N_s},
\quad
I_{\mathrm{bin}} = \sum_{n=0}^{N_s-1} i[n] \, e^{-j 2\pi n / N_s},
\label{eq:bin}
\end{equation}
\begin{equation}
H_{\mathrm{meas}}(j\omega_0) \;=\; I_{\mathrm{bin}} / V_{\mathrm{bin}}.
\label{eq:Hmeas}
\end{equation}
The cost surface for joint estimation of $(\alpha, R_x)$ is
\begin{equation}
J(\alpha, R_x) \;=\;
   \bigl| H_{\mathrm{meas}} - H(j\omega_0;\alpha, R_x) \bigr|^{2},
\label{eq:Jeuclid}
\end{equation}
with $H(j\omega_0;\cdot,\cdot)$ from (\ref{eq:abcd}).  The
maximum-likelihood variant under the proper-complex-Gaussian-ratio
noise model of \S\,\ref{sec:crlb} reweights the residual by
$\sigma_H^2(\alpha, R_x)$ from Appendix~B; both cost forms are
exercised in \S\,\ref{sec:results}.  Among the impedance / spectral
methods compared in \S\,\ref{sec:compare}, \cite{Nunes2019IJEPES}
is the closest single-ended antecedent; the v1 manuscript's [2] and
[10] both refer to this single entry (P0.3 de-duplication).

\section{Two-Stage Optimiser}\label{sec:opt}

The estimator solves $\min_{\theta} J(\theta)$ over
$\theta = (\alpha, R_x) \in (0,1)\times(\SI{1}{\ohm}, \SI{1}{\mega\ohm})$
via two stages.

\paragraph{Stage~1 (coarse search).}
A $100\!\times\!50$ grid is laid over the operating box, with the
$\alpha$ axis on a uniform grid and the $R_x$ axis on a
\emph{geometric} grid spanning six decades.  The geometric spacing
is essential: a uniform $R_x$ grid concentrates 95\,\% of the
sample budget above $R_x = \SI{1e5}{\ohm}$ where the cost surface
is flat (the high-impedance limit), starving the
$R_x \in [10^2, 10^4]\,\Omega$ regime where the
identifiability information lies.  The top three local minima of
$J$ on the grid seed Stage~2.

\paragraph{Stage~2 (continuous refinement).}
Each seed is refined by gradient descent with backtracking-Armijo
line search and box constraints.  The cost gradient is the
closed-form analytic expression from Appendix~A~\S\,A.7:
\begin{equation}
\frac{\partial J}{\partial \theta} \;=\;
   -2\,\mathrm{Re}\!\Bigl[
        \bigl(H_{\mathrm{meas}} - H(\theta)\bigr)^{*}\,
        \frac{\partial H(\theta)}{\partial \theta}
   \Bigr],
\label{eq:gradJ}
\end{equation}
with $\partial H/\partial\theta$ evaluated via the boxed quotient-rule
identity (\ref{eq:dHdtheta_quotient}).
The descent step direction is the Gauss--Newton diagonal
$p = -g \oslash h_{\mathrm{diag}}$ with $h_{\mathrm{diag}} =
2\,|\partial H/\partial \theta|^{2}$ (and a fallback to plain
$-g$ when $h_{\mathrm{diag}}$ is degenerate); this is the
analogue of one Newton iteration on a quadratic local model and
is the device that makes the Stage-2 iteration robust on cells
near the cost-surface degeneracy of \S\,\ref{sec:crlb}.

\paragraph{Analytical-vs-FD gradient swap (WP2.4).}
The analytical gradient (\ref{eq:gradJ}) replaces the central
finite-difference gradient used in v1 / Phase~1.  At convergence,
the two gradients yield the same fixed point on well-conditioned
cells (verified in \texttt{tests/test\_optimiser\_analytical\_vs\_fd
.py} to $|\Delta\hat\alpha| < 10^{-3}$,
$|\Delta\hat{R}_x|/R_x < 5\!\times\!10^{-3}$).  The finite-difference
gradient costs five $J$ evaluations per iteration (four for the two
partials plus one for the Armijo reference); the analytical gradient
costs zero extra evaluations beyond the line search.  Empirically the
WP2.4 swap reduces wall-clock per cell by 50--65\,\% at fixed
iteration cap, freeing budget for the 100-trial Monte-Carlo of
\S\,\ref{sec:results}.

\paragraph{Implementation.}
The two-stage estimator is implemented in
\texttt{inverse\_estimation/faultloc\_two\_stage\_optimiser.py}.
Defaults (Phase~2 onwards): \texttt{forward\_model='distributed'},
\texttt{gradient='analytical'}, \texttt{cost='ml'}.  The Phase-1
baseline (\texttt{forward\_model='cascaded\_gamma'},
\texttt{gradient='fd'}, \texttt{cost='euclid'}) is retained on the
backward-compatibility path for the §II-D legacy-comparison study.

\section{Simulation Setup and Noise Model}\label{sec:sim}
\textit{[\S\,V-A simulation programme, \S\,V-B dual-channel noise
model.  All headline numerical values cited here MUST reuse the
macros defined in the preamble - never hand-type
\samplingConfig{} or any of the headline percentages.]}

\section{Results and Discussion}\label{sec:results}

We report results in four layers, indexed against the project's
Phase~1 / Phase~2 acceptance keys (KPIs K01--K04 in the v3
Execution Plan~\S\,11).  Every numerical claim re-uses the
preamble macros so that the abstract, body, conclusion and
\texttt{docs/glossary.md} agree byte-identically.

% --- VI.A Phase-1 self-consistent estimator headline (unchanged) -----
\subsection{Self-consistent estimator headline (\S\,VI.A; KPI K01)}
\label{sec:resA}
On the 720-case grid sampled at \samplingConfig{} with waveforms
\emph{drawn from the optimiser's own forward model} (Phase~1
self-consistent path), the two-stage estimator achieves a mean
fault-location error of \headlineLocErrNoiseless{} in the noiseless
baseline and \headlineLocErrSNR{}\glossref{LOC-ERR}.  The mean
arc-resistance error is \headlineRxErrNoiseless{} noiseless and
\headlineRxErrEnvelope\glossref{RX-ERR}.  These numbers were
established under WP1.4 and are unchanged by the WP2.4 gradient
swap (the analytical gradient and the central-FD gradient reach the
same fixed point on the well-conditioned majority of the grid).

% --- VI.B Forward-model accuracy (K03) -------------------------------
\subsection{Forward-model accuracy (\S\,VI.B; KPI K03)}
\label{sec:resB}
Acceptance D-C of Phase~2 requires that the closed-form
distributed-parameter model (\ref{eq:abcd}) reproduce a 50-section
$\pi$ reference to better than $1\,\%$ in $|H|$ across the
operating grid.  The empirical envelope is \headlineKThreeAcc{};
the per-cell breakdown is in \texttt{outputs/phase2\_modelfit.csv}
(50 rows) and \texttt{outputs/phase2\_reproduction.csv} (95 rows on
the tightened $19\alpha\!\times\!5R_x$ grid of WP2.3).  K03
\textbf{passes with four orders of magnitude margin} --- the
residual is set by the 50-section reference's own discretisation
floor ($\sim 2\!\times\!10^{-5}\,\%$ vs a 100-section reference)
rather than a real distributed-vs-$\pi$ gap.  Table~\ref{tab:modelfit}
reports the head-to-head comparison against the v1 R-L-only and
P0.5 cascaded-$\Gamma$ baselines.  The 5$\times$ regression on the
single worst $(\alpha, R_x)=(0.95, 5000)$ cell of WP2.3
(\texttt{tests/test\_phase2\_no\_3944\_ceiling.py}) further certifies
\headlineFiveXReg{}; this is the Phase~2 retirement of the v1
manuscript's $39.44\,\%$ modelling-error ceiling.

% --- VI.C Cross-platform validation (Phase 1) ------------------------
\subsection{Cross-platform validation (\S\,VI.C; KPI K02)}
\label{sec:resC}
The estimator is exercised against three independent waveform
sources: PSCAD (\texttt{data/pscad\_720.mat}), EMTP-RV
(\texttt{data/emtp\_720.mat}) and a 50-section pure-MATLAB reference
(\texttt{data/ref\_50section\_720.mat}).  At noiseless operation the
mean cross-platform location error is $\sim 19\,\%$ (max $\sim
180\,\%$); at $\mathrm{SNR}_I \geq 30$\,dB it is $\sim 23$--$25\,\%$.
The 100-trial Monte-Carlo of WP1.5 (288\,000 trial-cells across the
four datasets) finds \emph{every} cell on the operating grid to be
statistically significantly biased ($p<0.05$, one-sided $t$-test
against zero), confirming that the residual is
\emph{structural} --- a single-bin DFT identifiability floor over a
near-degenerate curve in $(\alpha, R_x)$ --- rather than a
noise-realisation artefact.  The corresponding empirical CDF panels
are in \texttt{outputs/phase1\_figs/mc\_distribution\_*.png} and
the corrected-CRLB overlay is in
\texttt{outputs/phase1\_crlb\_overlay/}.

% --- VI.D Phase-2 cross-platform: K04 honest report ------------------
\subsection{Phase-2 cross-platform with continuously parametrised
forward model (\S\,VI.D; KPI K04)}\label{sec:resD}
We re-ran the WP1.4 cross-platform experiment using the
distributed-parameter forward model (\ref{eq:abcd}) and the
analytical-gradient optimiser of \S\,\ref{sec:opt}.  The K04
estimator-improvement metric --- mean Phase-2 vs Phase-1 cross-
platform reduction in location error at $\mathrm{SNR}_I \leq 30$\,dB
--- is \headlineKFourImp{}.  The per-cell delta is in
\texttt{outputs/phase2\_vs\_phase1\_delta.csv} and the per-trial
parquet is \texttt{outputs/phase2\_results\_per\_dataset.parquet};
six summary figures with the corrected-CRLB overlay are in
\texttt{outputs/phase2\_figs/}.

\textbf{Diagnosis.}  The forward-model fidelity has improved by
roughly five orders of magnitude (\S\,\ref{sec:resB}), but the
optimiser's location error is unchanged because the dominant source
of estimator error in the cross-platform regime is \emph{not} the
forward-model gap.  It is the cost-surface degeneracy of the single-
bin DFT identifiability floor (\S\,\ref{sec:resC}), which a forward-
model swap cannot break by construction: improving how accurately
$H(\alpha, R_x)$ is computed does not flatten or narrow the curve
in $(\alpha, R_x)$ along which $J(\alpha, R_x)$ is near-degenerate.
This finding is not a Phase~2 regression; it is the empirical
certification of an identifiability bound that the corrected
CRLB of \S\,\ref{sec:crlb} predicts.  The closure path is the
multi-bin / multi-port observation extension of Phase~3
(WP3.5 Taylor--Fourier estimator and WP3.6 multi-port FIM,
roadmap item~5 in \S\,\ref{sec:conclusion}) and is the subject of
the planned IEEE TPWRD follow-on submission.

% Headline-number echo block kept for the WP0.1 byte-identical check.
\subsection*{Headline numerical findings (WP0.1 echo)}
Phase-1 self-consistent envelope: location error
\headlineLocErrNoiseless{} noiseless and \headlineLocErrSNR{}
(Glossary \texttt{LOC-ERR}); $R_x$ error
\headlineRxErrNoiseless{} noiseless and \headlineRxErrEnvelope{}
(Glossary \texttt{RX-ERR}); empirical RMS \headlineCRLBGap{}
(Glossary \texttt{CRLB-GAP}).
Phase-2 forward-model accuracy: \headlineKThreeAcc{}.
Phase-2 cross-platform estimator improvement: \headlineKFourImp{}.
Phase-2 worst-cell regression: \headlineFiveXReg{}.
All experiments use \samplingConfig{} (Glossary \texttt{SAMP-CFG}).

% --- Figures (WP0.4 sub-task 6 - full captions, axis labels with units) ---
\begin{figure}[!t]
  \centering
  % Generated by matlab/figs/fig_section_convergence.m
  \fbox{\parbox[c][4.4cm][c]{0.92\linewidth}{\centering
    \textit{Figure placeholder} -- TF-magnitude modelling error
    of the $N_s$-section $\pi$-model versus a 50-section reference,
    swept over $N_s\in\{2,5,10,20\}$ and $\alpha\in[0.05,0.95]$ at
    $R_x=1\,000$\,$\Omega$, $f_0 = 50$\,Hz.\\
    (PDF: \texttt{outputs/fig\_section\_convergence.pdf})}}
  \caption{Section-count convergence study.  Horizontal axis:
           per-unit fault location $\alpha$ (dimensionless).
           Vertical axis: TF-magnitude modelling error vs the
           50-section reference, in percent.  $N_s = 2$ is the model
           used by the optimiser of \S\,\ref{sec:opt}; deeper
           sectionalisation is treated as the reference.  $R_x$ is
           held at $1\,000$\,$\Omega$ and the source frequency at
           $f_0 = 50$\,Hz.}
  \label{fig:section_convergence}
\end{figure}

\begin{figure}[!t]
  \centering
  % Generated by matlab/figs/fig_snr_sweep.m
  \fbox{\parbox[c][4.4cm][c]{0.92\linewidth}{\centering
    \textit{Figure placeholder} -- mean fault-location error vs
    current-channel SNR$_I$ at fixed voltage-channel SNR$_V$, averaged
    across the $720$-case grid.\\
    (PDF: \texttt{outputs/fig\_snr\_sweep.pdf})}}
  \caption{Mean fault-location error vs current-channel
           $\mathrm{SNR}_I$ (in dB; the rightmost cluster at 50\,dB on
           the abscissa denotes the noiseless case).  Vertical axis:
           mean per-cell relative location error in percent, on a
           logarithmic scale.  Each curve corresponds to a fixed
           $\mathrm{SNR}_V$ level
           ($\{20,30,40,\mathrm{noiseless}\}$\,dB).  The dominance of
           $\mathrm{SNR}_I$ over $\mathrm{SNR}_V$ at the same dB level
           is the empirical justification for the dual-channel noise
           model in \S\,\ref{sec:sim}.}
  \label{fig:snr_sweep}
\end{figure}

\begin{figure}[!t]
  \centering
  % Generated by matlab/figs/fig_alpha_rx_heatmap.m
  \fbox{\parbox[c][4.4cm][c]{0.92\linewidth}{\centering
    \textit{Figure placeholder} -- mean fault-location error
    heatmap over the $\alpha\times R_x$ plane, averaged over the
    $4\!\times\!4$ noise grid per cell.\\
    (PDF: \texttt{outputs/fig\_alpha\_rx\_heatmap.pdf})}}
  \caption{Mean fault-location error across the $9\,\alpha\!\times\!
           5\,R_x$ grid, averaged over the
           $4\,\mathrm{SNR}_V\!\times\!4\,\mathrm{SNR}_I$ noise
           classes.  Horizontal axis: arc resistance $R_x$ in
           ohms (logarithmic).  Vertical axis: per-unit fault
           location $\alpha$ (dimensionless).  Colour: $\log_{10}$ of
           the mean per-cell relative location error in percent;
           lighter shading indicates lower error.  Near-source
           behaviour ($\alpha < 0.2$) is reported separately in
           \S\,\ref{sec:crlb} as it is dominated by the CRLB floor
           rather than estimator behaviour.}
  \label{fig:alpha_rx_heatmap}
\end{figure}

\begin{figure}[!t]
  \centering
  % Generated by run_faultloc_phase2_continuous_param.py
  \includegraphics[width=0.92\linewidth]%
    {../outputs/phase2_figs/c_snrI_sweep_with_crlb.png}
  \caption{(Phase-2, WP2.5) Median per-cell location error vs
           current-channel $\mathrm{SNR}_I$ (dB) on the three
           cross-platform datasets (PSCAD, EMTP-RV, 50-section
           reference), with the proper-complex-Gaussian-ratio CRLB
           (dashed) and the dual-channel CRLB (dotted) of
           \S\,\ref{sec:crlb} overlaid.  The optimiser uses the
           continuously parametrised forward model
           (\ref{eq:abcd}) and the analytical-gradient
           (\ref{eq:gradJ}) of \S\,\ref{sec:opt}.  The empirical
           curves \emph{do not} converge to the CRLB envelope below
           $\mathrm{SNR}_I = 30$\,dB; the residual is the single-bin
           DFT identifiability floor of \S\,\ref{sec:resC} and
           closes at WP3.5 / WP3.6 (Taylor--Fourier multi-bin +
           multi-port FIM).}
  \label{fig:phase2_snrI_crlb}
\end{figure}

\section{Comparison with Existing Methods}\label{sec:compare}
\textit{[\S\,VII Table 3-bis with four numerical competitors lands at
WP4.5 (Phase 4): Paramo-2023 eigenvalue, Iurinic-2018 spectral,
Cui-Weng-2020 micro-PMU, Zeng-2021 double-ended HIF.]}

\section{Cram\'er--Rao Lower Bound Analysis}\label{sec:crlb}

The bound used by the v1 manuscript treats $H_{\mathrm{meas}}$ as
contaminated by circular complex Gaussian noise with constant
variance.  This linearisation is \emph{not} self-consistent with the
\S\,V-B dual-channel AWGN noise model: $H_{\mathrm{meas}} =
I_{\mathrm{bin}}/V_{\mathrm{bin}}$ is the ratio of two independent
complex Gaussians, whose density is the
Marsaglia--Nadarajah--Pog\'any form~\cite{Marsaglia1965Ratio,
NadarajahPogany2018,Kuruoglu2018ASCE}, not Gaussian.  The Gaussian-on-$H$
linearisation is valid only in the asymptotic
$|I| \gg \sigma_I$ regime --- the opposite of the HIF regime,
where $|H|^2 \sigma_V^2$ becomes comparable to $\sigma_I^2$.

We therefore derive two corrected bounds (full derivation in
Appendix~B):

\paragraph{Proper-complex-Gaussian-ratio FIM.}
At high SNR$_V$ (Geary--Hinkley~\cite{Marsaglia1965Ratio}
$|V_{\mathrm{phase}}|/\sigma_V > 4$, satisfied across the entire
720-case grid for time-domain SNR$_V > -8$\,dB) the ratio density
admits a complex-Gaussian approximation with a
\emph{signal-dependent} per-component variance
$\sigma_H^2 = (\sigma_I^2 + |H|^2 \sigma_V^2)/|V_{\mathrm{phase}}|^2$.
The resulting FIM is
$F^{\mathrm{proper}}_{kl} = (1/\sigma_H^2)\,
\mathrm{Re}[(\partial_{\theta_k} H)^* (\partial_{\theta_l} H)]$,
implemented in \texttt{inverse\_estimation/faultloc\_crlb\_proper.py}.

\paragraph{Joint dual-channel FIM (Nehorai--Hawkes, 2000)~\cite{NehoraiHawkes2000}.}
Working directly in $(V(t), I(t))$ time-domain space and projecting
onto $(\alpha, R_x)$ via the chain rule with $\partial V/\partial\theta
\equiv 0$ and $\partial I/\partial\theta = \mathrm{Re}[(\partial H/
\partial \theta) V_{\mathrm{phase}} e^{j\omega_0 t}]$ gives the
ground-truth bound in waveform space, implemented in
\texttt{inverse\_estimation/faultloc\_crlb\_dualchannel.py}.

\paragraph{Cross-check.}
The two FIMs share a common Fisher kernel and differ in their
prefactor; their ratio is $F^{\mathrm{proper}}/F^{\mathrm{dual}} =
\sigma_I^2/(\sigma_I^2 + |H|^2 \sigma_V^2)$.  At $\sigma_V \to 0$
(V perfectly known) the two are exactly equal; at finite $\sigma_V$
the proper-ratio bound is looser by exactly the factor expected from
the loss of information when only the H-ratio is observed.  Both
bounds are tested for consistency in
\texttt{tests/test\_crlb\_consistency.py} (5 tests, all pass).

\paragraph{Four headline findings (WP1.6 corrected).}
\begin{enumerate}
  \item The proper-ratio bound supersedes the v1 Gaussian-on-$H$
        linearisation; the latter under-estimates the true bound and
        is invalid in the HIF regime.
  \item The dual-channel bound is tighter than the proper-ratio bound
        whenever $\sigma_V > 0$, by a factor that quantifies the
        information loss from working in H-ratio space rather than
        on the raw waveforms.
  \item The empirical RMS estimator stays \headlineCRLBGap{} of the
        proper-ratio CRLB across the high-SNR cells of the 720-case
        grid (Glossary entry \texttt{CRLB-GAP}; overlay panel
        \texttt{outputs/phase1\_crlb\_overlay/}).
  \item The Geary--Hinkley validity flag is satisfied across the
        full 720-case grid; the Gaussian-on-ratio approximation is
        therefore appropriate for the WP1.6 reporting envelope.
\end{enumerate}

% =============================================================================
\section{Conclusion and Future Work}\label{sec:conclusion}
% =============================================================================
This paper presented a single-ended joint estimator of HIF location and
arc resistance from a single power-frequency admittance bin under a
dual-channel additive-white-Gaussian-noise model.  Numerically, the
estimator delivers a mean fault-location error of
\headlineLocErrNoiseless{} (noiseless) and \headlineLocErrSNR{} on
the self-consistent path of \S\,VI-A.  The arc-resistance envelope
is the principal scope statement of the estimator:
\headlineRxErrEnvelope.  The empirical RMS error remains
\headlineCRLBGap.  The Phase-2 upgrade (continuously parametrised
forward model with closed-form analytical gradient) attains
\headlineKThreeAcc{} forward-model fidelity vs the 50-section
$\pi$ reference and \headlineFiveXReg{} on the v1 ceiling cell;
the cross-platform estimator-improvement metric (KPI K04) is
\headlineKFourImp{}, a result honestly reported in \S\,VI-D and
attributed to the single-bin DFT identifiability floor of
\S\,VIII --- a structural limit the corrected CRLB predicts and
that the planned multi-bin / multi-port observation extension
breaks.

\subsection{Roadmap (12 / 24 / 36 months)}
The future work below is anchored to Phases 1--5 of the SAMBPS DTaaS
\emph{HIF-TF Project Execution Plan v1.0}
(\texttt{docs/FaultLocationIdentification\_ExecutionPlan.pdf},
\S\,2 and \S\,4).

\noindent\textbf{12-month horizon (Phases 1--2, decision gates D1--D2).}
\begin{enumerate}
  \item Cross-platform validation on PSCAD, EMTP-RV and a 50-section
        pure-MATLAB reference, plus 100-trial Monte-Carlo per grid
        cell with 95th-percentile confidence intervals
        (Phase 1, WP1.1--WP1.5).
  \item Re-derivation of the Fisher Information Matrix under the
        proper-complex-Gaussian-ratio framework
        (Kuruo\u{g}lu-2018) and the joint dual-channel FIM
        in $(V,I)$ space (Nehorai-Hawkes-2000), with the
        Geary-Hinkley validity test reported per cell
        (Phase 1, WP1.6).
  \item Continuously parametrised distributed-parameter $H(j\omega_0;
        \alpha, R_x)$ via cascaded ABCD / Bergeron blocks, retiring
        the 39.44\,\% two-section modelling-error ceiling, with
        analytical $\partial H/\partial\alpha,\partial H/\partial R_x$
        replacing the central-difference gradient
        (Phase 2, WP2.1--WP2.5).
\end{enumerate}

\noindent\textbf{24-month horizon (Phases 3--4, decision gates D3--D4).}
\begin{enumerate}\setcounter{enumi}{3}
  \item Generalisation to three-phase $Y_{abc}(j\omega_0)$ on the IEEE
        13-, 34- and 123-node test feeders with laterals, tap loads
        and at least one distributed generator
        (Phase 3, WP3.1--WP3.4).
  \item First-order Taylor--Fourier phasor estimator
        (Platas-Garza-2010) replacing the single-bin DFT, plus
        Hermann--Krener / STRIKE-GOLDD identifiability mapping of
        $J(\alpha, R_x)$, and a multi-port FIM extension
        (Phase 3, WP3.5--WP3.6).
  \item Field-grade impairment classes (impulsive, harmonic,
        CT/PT saturation per IEEE C37.110, $\pm 0.5$\,Hz off-nominal
        per IEEE C37.118.1 P-class, ADC quantisation) and arc-class
        diversification (Cassie--Mayr--Kizilcay, Wang-2020
        distortion-controllable, Torres-2022 stochastic), plus a
        numerical Table 3-bis benchmarking four competitor methods
        on identical 720-case waveforms
        (Phase 4, WP4.1--WP4.5).
\end{enumerate}

\noindent\textbf{36-month horizon (Phase 5, decision gate D5).}
\begin{enumerate}\setcounter{enumi}{6}
  \item RTDS / Typhoon HIL pilot at IIT-Madras with real relay-class
        IED integration via IEC 61850-9-2 SV; live HIF arc scenarios
        targeting end-to-end estimate latency below five
        power-frequency cycles (Phase 5, WP5.1--WP5.3).
  \item Productisation as the SAMBPS DTaaS Protection-Validation
        module v1.0 with REST + Python client API and identifiability
        + CRLB overlays surfaced to the operator
        (Phase 5, WP5.4).
  \item Second journal paper (IEEE TPWRD or TSG) bringing in HIL
        evidence, the four-competitor numerical benchmark and the
        multi-port CRLB (Phase 5, WP5.5).
\end{enumerate}

\subsection{Integration into the SAMBPS DTaaS Protection-Validation
            module}
The estimator's single-frequency, training-free, single-ended
operation maps cleanly onto the SAMBPS DTaaS scenario-engine: the
digital twin streams synthetic dual-channel V/I waveforms through the
locator, surfaces the identifiability flag from
\texttt{adaptation/faultloc\_identifiability\_check.py}, and overlays
the corrected CRLB envelope on the operator's location estimate.
This binds the estimator to a maintainable validation harness rather
than a one-shot study, and forms the deployment vehicle for the v1.0
release at decision gate D5
(see Execution Plan \S\,4.6 WP5.4 and \S\,11 KPI K16).

% =============================================================================
\section*{Acknowledgments}
% =============================================================================
The authors thank the IIT-Madras Power Systems Computational Lab for
hosting the experimental work and SAMBPS Digital Twin Labs for
sponsoring the reproducibility-grade software release of all PSCAD
and EMTP-RV models, the 720-case waveform set and the 50-section
reference dataset under MIT licence.

% =============================================================================
\bibliographystyle{IEEEtran}
\bibliography{references}
% =============================================================================

\end{document}
